\section{Proposed Graph-VR-Det Methodology}

\subsection{Overall Architecture}
The Graph-VR-Det framework reformulates the near-field XL-MIMO detection task as a regression problem on a dynamically constructed sparse bipartite graph. Let $\mathcal{G} = (\mathcal{V}, \mathcal{U}, \mathcal{E})$ denote an undirected bipartite graph, where $\mathcal{V} = \{v_1, v_2, \dots, v_N\}$ represents the set of $N$ base station antenna nodes, and $\mathcal{U} = \{u_1, u_2, \dots, u_K\}$ represents the set of $K$ user nodes. The edge set $\mathcal{E}$ defines the connectivity between antennas and users. In conventional fully-connected GNN-based detectors, $\mathcal{E}$ contains $NK$ edges, implying that every user's signal is processed by the entire array. However, due to the spatial non-stationarity inherent in near-field XL-MIMO channels, the signal energy of each user is concentrated within a localized visibility region (VR). Consequently, processing the full graph introduces severe noise accumulation from non-visible antennas and incurs prohibitive computational overhead. To address this limitation, Graph-VR-Det introduces a dynamic VR discovery mechanism to construct a sparse edge set $\mathcal{E}_{\text{sparse}} \subset \mathcal{E}$, where $|\mathcal{E}_{\text{sparse}}| = SK \ll NK$, with $S$ denoting the number of retained antennas per user. A sparse message passing neural network (Sparse-GNN) is then executed exclusively over $\mathcal{E}_{\text{sparse}}$ to iteratively refine the estimates of the transmitted symbols.

\subsection{Dynamic Visibility Region Discovery}
The construction of the sparse graph topology relies on a pilot-based dynamic VR discovery algorithm, operating without requiring perfect prior knowledge of the channel state information (CSI) or the true VR mask $\mathbf{M}$. During the uplink training phase, orthogonal pilot sequences are transmitted by the $K$ users. Let $\mathbf{X}_p = \mathbf{I}_K$ denote the pilot matrix. The received pilot signal at the BS is given by $\mathbf{Y}_p = \mathbf{H} + \mathbf{N}_p$, where $\mathbf{N}_p$ is the additive noise matrix during the pilot phase. To identify the effective VR for each user, the receiver computes the energy of the received pilot signals. The energy matrix $\mathbf{E} \in \mathbb{R}^{N \times K}$ is calculated element-wise as $E_{n,k} = |[\mathbf{Y}_p]_{n,k}|^2$. For the $k$-th user, the algorithm selects the indices of the antennas corresponding to the $S$ largest energy values in the $k$-th column of $\mathbf{E}$. This top-$S$ selection dynamically generates a sparse binary adjacency matrix $\mathbf{A}_{\text{sparse}} \in \{0, 1\}^{N \times K}$, where $[\mathbf{A}_{\text{sparse}}]_{n,k} = 1$ if the $n$-th antenna is within the discovered VR of the $k$-th user, and $0$ otherwise. The non-zero entries of $\mathbf{A}_{\text{sparse}}$ define the sparse edge set $\mathcal{E}_{\text{sparse}}$. By relying on the received signal energy rather than exact channel coefficients, this approach is robust to estimation errors and effectively filters out antennas dominated by noise.

\subsection{Sparse Message Passing Neural Network}
Given the sparse graph topology $\mathcal{E}_{\text{sparse}}$, the Sparse-GNN performs iterative evidence processing to recover the transmitted symbols. The process begins with the initialization of node and edge features. The initial feature vector for the $n$-th antenna node is defined by the real and imaginary parts of the received signal, i.e., $\mathbf{v}_n^{(0)} = [\Re(y_n), \Im(y_n)]^T \in \mathbb{R}^2$. The initial feature vector for the $k$-th user node is set to zero, $\mathbf{u}_k^{(0)} = [0, 0]^T \in \mathbb{R}^2$. The edge feature between the $n$-th antenna and the $k$-th user is initialized using the estimated channel coefficient, $\mathbf{e}_{n,k} = [\Re(H_{n,k}), \Im(H_{n,k})]^T \in \mathbb{R}^2$. These initial features are mapped to a higher-dimensional hidden space $\mathbb{R}^{d_h}$ via linear embedding layers, where $d_h$ is the hidden dimension.

The core of the Sparse-GNN consists of $L$ message passing layers. In the $l$-th layer, where $l \in \{1, 2, \dots, L\}$, the network performs Antenna-to-User (A2U) and User-to-Antenna (U2A) updates strictly along the edges in $\mathcal{E}_{\text{sparse}}$. During the A2U phase, the message from antenna $n$ to user $k$ is computed using a multi-layer perceptron (MLP), denoted as $\text{MLP}_{\text{A2U}}$, which takes the concatenated features of the connected nodes and the edge as input:
\begin{equation}
    \mathbf{m}_{n \to k}^{(l)} = \text{MLP}_{\text{A2U}} \left( \left[ \mathbf{v}_n^{(l-1)}, \mathbf{u}_k^{(l-1)}, \mathbf{e}_{n,k} \right] \right).
\end{equation}
The messages are then aggregated at the $k$-th user node using a scatter-add operation over the valid edges, and the user feature is updated via $\text{MLP}_{\text{U}}$:
\begin{equation}
    \mathbf{u}_k^{(l)} = \text{MLP}_{\text{U}} \left( \left[ \mathbf{u}_k^{(l-1)}, \frac{1}{S} \sum_{n \in \mathcal{N}(k)} \mathbf{m}_{n \to k}^{(l)} \right] \right),
\end{equation}
where $\mathcal{N}(k)$ denotes the set of antennas connected to user $k$ in $\mathcal{E}_{\text{sparse}}$. Subsequently, during the U2A phase, messages from users to antennas are computed and aggregated similarly:
\begin{equation}
    \mathbf{m}_{k \to n}^{(l)} = \text{MLP}_{\text{U2A}} \left( \left[ \mathbf{u}_k^{(l)}, \mathbf{v}_n^{(l-1)}, \mathbf{e}_{n,k} \right] \right),
\end{equation}
\begin{equation}
    \mathbf{v}_n^{(l)} = \text{MLP}_{\text{A}} \left( \left[ \mathbf{v}_n^{(l-1)}, \frac{1}{|\mathcal{N}(n)|} \sum_{k \in \mathcal{N}(n)} \mathbf{m}_{k \to n}^{(l)} \right] \right),
\end{equation}
where $\mathcal{N}(n)$ is the set of users connected to antenna $n$, and $|\mathcal{N}(n)|$ is its dynamic degree. After $L$ iterations, a linear readout layer maps the final user features $\mathbf{u}_k^{(L)}$ to a two-dimensional space to predict the real and imaginary parts of the transmitted symbols $\hat{\mathbf{x}}$.

\subsection{Frequency-Aware Wideband Extension}
To extend the architecture to wideband THz XL-MIMO systems, the near-field beam splitting effect must be addressed. Beam splitting causes the spatial focusing point, and consequently the VR, to shift across different subcarriers. Processing each subcarrier independently with a separate graph topology is computationally inefficient and ignores the underlying physical correlation of the wideband channel. Therefore, Graph-VR-Det introduces a joint frequency-domain graph mapping strategy. During the pilot phase, the energy matrix is computed for each subcarrier $m$, denoted as $\mathbf{E}_m$. The algorithm then calculates the average energy across all $F$ subcarriers, $\bar{\mathbf{E}} = \frac{1}{F} \sum_{m=1}^F \mathbf{E}_m$. The top-$S$ selection is performed on $\bar{\mathbf{E}}$ to generate a shared, stable graph topology $\mathcal{E}_{\text{sparse}}$ that is applied universally across all subcarriers.

To enable the GNN to learn and compensate for the frequency-dependent spatial shifts within this shared topology, frequency-aware edge features are introduced. For the $m$-th subcarrier, the normalized frequency offset $\Delta f_m = (f_m - f_c) / f_c$ is concatenated with the channel coefficients to form an augmented three-dimensional edge feature vector:
\begin{equation}
    \mathbf{e}_{n,k,m} = \left[ \Re(H_{n,k,m}), \Im(H_{n,k,m}), \frac{f_m - f_c}{f_c} \right]^T \in \mathbb{R}^3.
\end{equation}
By embedding the explicit frequency information into the message passing kernel, the Sparse-GNN dynamically adjusts its evidence processing weights to mitigate the beam splitting effect, effectively fusing multi-carrier spatial features without requiring subcarrier-specific network parameters.

\subsection{Complexity Analysis}
The computational complexity of Graph-VR-Det is significantly lower than that of traditional linear detectors and fully-connected GNNs. For full-array MMSE detection, the complexity is dominated by the matrix inversion $(\mathbf{H}^H \mathbf{H} + \sigma^2 \mathbf{I}_K)^{-1}$ and the matrix multiplication $\mathbf{H}^H \mathbf{H}$, which requires $\mathcal{O}(K^3 + NK^2)$ floating-point operations (FLOPs) per subcarrier. As the number of antennas $N$ scales to the thousands in XL-MIMO systems, this cubic/quadratic scaling becomes a severe bottleneck. For a fully-connected MPNN (Full-MPNN), the message passing operations are executed over all $NK$ edges. With $L$ layers and a hidden dimension $d_h$, the complexity scales as $\mathcal{O}(L N K d_h^2)$, which still grows linearly with the total number of antennas $N$, leading to potential out-of-memory (OOM) issues.

In contrast, the Sparse-GNN restricts all message passing operations strictly to the $SK$ edges defined by $\mathcal{E}_{\text{sparse}}$. The complexity of the dynamic VR discovery is $\mathcal{O}(NK \log S)$ for the top-$S$ sorting. The subsequent message passing complexity is reduced to $\mathcal{O}(L S K d_h^2)$. Since the VR size $S$ is determined by the physical array aperture and user distance, it remains relatively constant and satisfies $S \ll N$. Consequently, the computational complexity of the Sparse-GNN scales as $\mathcal{O}(SK)$, effectively decoupling the detection complexity from the massive total number of antennas $N$. This reduction from $\mathcal{O}(N K)$ to $\mathcal{O}(S K)$ not only minimizes the memory footprint but also enables zero-shot scalability to extremely large arrays without retraining, as the local graph structure within the VR remains consistent regardless of the global array size.